% Double stroked N for natural numbers
\newcommand{\N}{\ensuremath{\mathbb{N}}}

% Double stroked Z for integers
\newcommand{\Z}{\ensuremath{\mathbb{Z}}}

% Double stroked Q for rational numbers
\newcommand{\Q}{\ensuremath{\mathbb{Q}}}

% Double stroked R for reals
\newcommand{\R}{\ensuremath{\mathbb{R}}}

% Double stroked K for a field (K as in "`Körper"'). Very common in Banach 
% space analysis. Same for F (as in field) for texts in english.
\newcommand{\K}{\ensuremath{\mathbb{K}}}
\newcommand{\F}{\ensuremath{\mathbb{F}}}

% Double stroked E, commonly used in discrete stochastics and asymptotic
% functional analysis to denote the expected value

\newcommand{\E}{\ensuremath{\mathbb{E}}}

% Double stroked P for the set of prime numbers
\newcommand{\Prim}{\ensuremath{\mathbb{P}}}

% Double stroked 1. Commonly used to denote indicator functions.
\newcommand{\Eins}{\ensuremath{\mathds{1}}}

% For flooring an expression without opening and closing any tags
\newcommand{\floor}[1]{\left\lfloor{#1}\right\rfloor}

% For ceiling an expression without opening and closing any tags
\newcommand{\ceil}[1]{\left\lceil{#1}\right\rceil}

% The following four commands create a quick and dirty way to type a sum,
% product, cap or cup in a $.$-math environment having the bounds above and 
% below the sign instead of its right. Use them economically as they
% destroy any flow.

\newcommand{\Sum}[2]{\stackrel{#2}{\underset{#1}{\sum}}}
\newcommand{\Prod}[2]{\stackrel{#2}{\underset{#1}{\prod}}}
\newcommand{\altCup}[2]{\overset{#2}{\underset{#1}{\bigcup}}}
\newcommand{\altCap}[2]{\overset{#2}{\underset{#1}{\bigcap}}}
\newcommand{\Sup}[1]{\underset{#1}\sup\,}
\newcommand{\Inf}[1]{\underset{#1}\inf\,}

% Simplified command for lower/upper integral bounds
\newcommand{\Int}[2]{\int\limits_{#1}^{#2}}

% Simplified command for an integral with just one lower bound, as in curvature
% integrals or integrals in measure theory
\newcommand{\Intmen}[1]{\underset{#1}{\int}}

% These three commands create brackets around the variable and add its 
% (single!) index (subscript) and its range. Handle with care ;).
% Requires the altstyle.sty to create the "Laue-Haken".

\newcommand{\finfol}[3]{(#1_{#2})_{#2\in\haken{#3}}}   %finite series
\newcommand{\finfolN}[3]{(#1_{#2})_{#2\in\haken{#3}_0}}%finite series with zero
\newcommand{\folge}[3]{(#1_{#2})_{#2 \in #3}}          %series

% Denotes the span of the argument in algebraic notation
\newcommand{\erz}[1]{\left\langle #1\right\rangle}

% Denotes the absolute value of the argument
\newcommand{\abs}[1]{\left|#1\right|}

% Use this to create "sets". The first argument defines the variable (x \in X
% for instance), the second one its defining property (e. g. x \neq y).

\newcommand{\menge}[2]{\left\{#1\,|\,#2\right\}}
\newcommand{\mengea}[2]{\left\{#1\,\vrule\,#2\right\}}

% Puts set brackets around the argument

\newcommand{\meng}[1]{\left\{#1\right\}}

% Same as \abs, but with "norm lines".
\newcommand{\norm}[1]{\lVert#1\rVert} %simple norm, always small delimiters
\newcommand{\norma}[1]{\left\|#1\right\|} %variable delimiters

% Creates the commonly used notation of the "operator norm". The first
% argument is the operator, the second its domain and the third its range.

\newcommand{\opnorm}[3]{\left\|#1\right\|_{\scriptscriptstyle{#2 \rightarrow #3}}}

% Puts [] around the argument
\newcommand{\klasse}[1]{\left[#1\right]}

% Simplifies the \lim command
\newcommand{\Limes}[2]{\underset{#1\to #2}{\lim}}

% limsup
\newcommand{\Lims}[2]{\underset{#1\to #2}{\overline{\lim}}}

% liminf
\newcommand{\Limi}[2]{\underset{#1\to #2}{\underline{\lim}}}

% lim from below
\newcommand{\Limup}[2]{\underset{#1\nearrow #2}{\lim}}

% lim from above
\newcommand{\Limdown}[2]{\underset{#1\searrow #2}{\lim}}

% creates a paragraph
\newcommand{\absatz}{\hfill \medskip \newline}

% creates a indented paragraph
\newcommand{\parag}{\hfill \medskip \par}

% denotes the "id" (as in identity)
\newcommand{\id}{\text{id}\,}

% Abbreviations
\newcommand{\Def}{\mathop{\rm Def}\nolimits}    % domain

\newcommand{\Kern}   {\mathop{\rm Kern}\nolimits}   % kernel
\newcommand{\Bild}   {\mathop{\rm Bild}\nolimits}   % range
\newcommand{\Spur}   {\mathop{\rm Spur}\nolimits}   % trace
\newcommand{\Inn}    {\mathop{\rm int}\nolimits}    % interior
\newcommand{\Syl}    {\mathop{\rm Syl}\nolimits}    % sylow groups
\newcommand{\minpol} {\mathop{\rm minPol}\nolimits} % minimal polynomial
\newcommand{\charpol}{\mathop{\rm charPol}\nolimits}% characteristic polynomial
\newcommand{\chark}  {\mathop{\rm char}\nolimits}   % characteristic (rings)
\newcommand{\grad}   {\mathop{\rm grad}\nolimits}   % degree
\newcommand{\sgn}    {\mathop{\rm sgn}\nolimits}    % sign
\newcommand{\Pot}    {\mathop{\rm Pot}\nolimits}    % power set
\newcommand{\End}    {\mathop{\rm End}\nolimits}    % endomorphisms
\newcommand{\Aut}    {\mathop{\rm Aut}\nolimits}    % automorphisms
\newcommand{\Hom}    {\mathop{\rm Hom}\nolimits}    % homomorphisms
\newcommand{\Epi}    {\mathop{\rm Epi}\nolimits}    % epimorphisms
\newcommand{\Mon}    {\mathop{\rm Mon}\nolimits}    % monomorphisms
\newcommand{\Iso}    {\mathop{\rm Iso}\nolimits}    % isomorphisms
\newcommand{\Fix}    {\mathop{\rm Fix}\nolimits}    % fix point set
\newcommand{\Inv}    {\mathop{\rm Inv}\nolimits}    % units
\newcommand{\ggt}    {\mathop{\rm ggT}\nolimits}    % greatest common divisor
\newcommand{\kgv}    {\mathop{\rm kgV}\nolimits}    % least common multiple
\newcommand{\inv}    {^{-1}}                        % to the power of -1
\newcommand{\ord}    {\mathop{\rm ord}\nolimits}    % degree of a group element
\newcommand{\Gal}    {\mathop{\rm Gal}\nolimits}    % galois group
\newcommand{\Mat}    {\mathop{\rm Mat}\nolimits}
\newcommand{\Umg}    {\mathop{\rm Umg}\nolimits}
\newcommand{\pr}     {\mathop{\rm pr}\nolimits}
\newcommand{\spann}  {\mathop{\rm span}\nolimits}
\newcommand{\conv}   {\mathop{\rm conv}\nolimits}
%induction stuff
\newcommand{\ia}{\textsl{Induktionsanker:}}
\newcommand{\iv}{\textsl{Induktionsannahme:}}
\newcommand{\is}{\textsl{Induktionsschritt:}}

% "`=>"' and "`<="'
\newcommand{\rueck}{"`$\Leftarrow$"':}
\newcommand{\hin}{"`$\Rightarrow$"':}

% same for subset and supset
\newcommand{\teilmenge}{"`$\subseteq$"':\\}
\newcommand{\obermenge}{"`$\supseteq$"':\\}

% some common distributions (not common notation)
\newcommand{\hyp}{\textsl{Hyp}} %hypergeometric
\newcommand{\ber}{\textsl{Ber}} %Bernoulli
\newcommand{\cau}{\textsl{Cau}} %Cauchy
\newcommand{\bin}{\textsl{Bin}} %binomial
\newcommand{\poi}{\textsl{Poi}} %Poisson
\newcommand{\Var}{\textsl{Var}} %variance

\newcommand{\erf}{\mathop{\rm erf}\nolimits}

% direct and orthogonal sum - if you find a single sign for this one, let
% me know ;)
\newcommand{\orth}{\text{\textcircled{$\bot$}}}

% litte shorter for single letters
\newcommand{\mca}{\mathcal{A}}
\newcommand{\mcb}{\mathcal{B}}
\newcommand{\mcc}{\mathcal{C}}
\newcommand{\mcd}{\mathcal{D}}
\newcommand{\mce}{\mathcal{E}}
\newcommand{\mcf}{\mathcal{F}}
\newcommand{\mcg}{\mathcal{G}}
\newcommand{\mch}{\mathcal{H}}
\newcommand{\mci}{\mathcal{I}}
\newcommand{\mcj}{\mathcal{J}}
\newcommand{\mck}{\mathcal{K}}
\newcommand{\mcl}{\mathcal{L}}
\newcommand{\mcm}{\mathcal{M}}
\newcommand{\mcn}{\mathcal{N}}
\newcommand{\mco}{\mathcal{O}}
\newcommand{\mcp}{\mathcal{P}}
\newcommand{\mcq}{\mathcal{Q}}
\newcommand{\mcr}{\mathcal{R}}
\newcommand{\mcs}{\mathcal{S}}
\newcommand{\mct}{\mathcal{T}}
\newcommand{\mcu}{\mathcal{U}}
\newcommand{\mcv}{\mathcal{V}}
\newcommand{\mcw}{\mathcal{W}}
\newcommand{\mcx}{\mathcal{X}}
\newcommand{\mcy}{\mathcal{Y}}
\newcommand{\mcz}{\mathcal{Z}}

\newcommand{\mfa}{\mathfrak{A}}
\newcommand{\mfb}{\mathfrak{B}}
\newcommand{\mfc}{\mathfrak{C}}
\newcommand{\mfd}{\mathfrak{D}}
\newcommand{\mfe}{\mathfrak{E}}
\newcommand{\mff}{\mathfrak{F}}
\newcommand{\mfg}{\mathfrak{G}}
\newcommand{\mfh}{\mathfrak{H}}
\newcommand{\mfi}{\mathfrak{I}}
\newcommand{\mfj}{\mathfrak{J}}
\newcommand{\mfk}{\mathfrak{K}}
\newcommand{\mfl}{\mathfrak{L}}
\newcommand{\mfm}{\mathfrak{M}}
\newcommand{\mfn}{\mathfrak{N}}
\newcommand{\mfo}{\mathfrak{O}}
\newcommand{\mfp}{\mathfrak{P}}
\newcommand{\mfq}{\mathfrak{Q}}
\newcommand{\mfr}{\mathfrak{R}}
\newcommand{\mfs}{\mathfrak{S}}
\newcommand{\mft}{\mathfrak{T}}
\newcommand{\mfu}{\mathfrak{U}}
\newcommand{\mfv}{\mathfrak{V}}
\newcommand{\mfw}{\mathfrak{W}}
\newcommand{\mfx}{\mathfrak{X}}
\newcommand{\mfy}{\mathfrak{Y}}
\newcommand{\mfz}{\mathfrak{Z}}

% common in probability theory and function theory
\newcommand{\spi}{\frac{1}{\sqrt{2\pi}}}

% \varepsilon looks great, but is way to long, same for \varphi
\newcommand{\eps}{\varepsilon}
\newcommand{\vi}{\varphi}

% tautology (quick and dirty)
\newcommand{\tteq}{\dashv \vdash}

% implication with extra space
\newcommand{\folgt}{\, \Rightarrow \,}

% creates two links to the same object, respecting the current value
% (ref). For instance "\shortref{fig:figure}{Fig.}"' will create a blue
% "`Fig. 1"', where both parts link to the figure.
\newcommand{\shortref}[2]{\hyperref[#1]{#2} \ref{#1}}

% inner produkt in "langle"-notation.
\newcommand{\sprod}[2]{\left\langle #1 |#2\right\rangle}

% used for the linear combinator (by me ;))
\newcommand{\bprod}[2]{\left\langle #1 ||#2\right\rangle}

% add the current section to the table of contents.

\newcommand{\tocsec}[2]
 {\addcontentsline{toc}{section}{\protect\numberline{#1} #2}}

% Some ball abbreviations.

% open ball in #1 with centre #2 and radius #3
\newcommand{\koff}[3]{K_{#1}\left(#2, #3\right)}

% closed ball in #1 with centre #2 and radius #3
\newcommand{\kabg}[3]{\overline K_{#1}\left(#2, #3\right)}

% the boundary of a closed ball in #1 with centre #2 and radius #3
\newcommand{\kran}[3]{\partial \kabg{#1}{#2}{#3}}

\newcommand{\koffk}[3]{K_{\scriptscriptstyle{#1}}
                        ^{\scriptscriptstyle{\left(#2, #3\right)}}}

% closed ball in #1 with centre #2 and radius #3
\newcommand{\kabgk}[3]{\overline K_{\scriptscriptstyle{#1}}
                        ^{\scriptscriptstyle{\left(#2, #3\right)}}}

% the boundary of a closed ball in #1 with centre #2 and radius #3
\newcommand{\krank}[3]{\partial \kabgk{#1}{#2}{#3}}

% overset an arrow with an index
\newcommand{\pfeil}[2]{ {\overrightarrow{#1}}^{{}_#2} }
\newcommand{\pfeilk}[2]{ {\scriptstyle{\overrightarrow{#1}}^{{}_#2}} }

\newcommand{\restrict}[2]{#1\vrule_{{}_{\scriptstyle #2}}}

\newcommand{\refclaim}[2]{\ref{#1}.(\ref{#2})}
\newcommand{\refcl}[2]{\ref{#1}.({\ref{#1#2}})}

\newcommand{\claim}[1]{\item \label{\thesection{#1}}}
